\section{軟弱地面に対する従来歩行制御手法の限界}

床沈み込みを伴う環境においては，従来の歩行制御手法の前提条件が部分的に破綻する．多くの手法では，床変形による影響を外乱として扱い，フィードバック制御によって補償することを想定している．しかし，床沈み込みは，接地と同時に発生し，かつ継続的に変化する現象であるため，単純な外乱補償では対応が困難である．

また，CoMの水平方向運動の制御のみに着目した手法では，CoM高さ変動に起因するZMP変動を直接的に抑制することができない．その結果，安定化制御が過度に作用し，運動が不自然になる，あるいは制御限界に達する可能性がある．

これらの理由から，軟弱地面における歩行では，従来の剛体床を前提とした歩行制御手法をそのまま適用することには限界がある．
